% Pursuing Stacks (cotes 134-2 to 134-8): archive cote 134-5, site page 83 (PDF page 84), typescript page 270.
% typescript. Draft compared: ps4.tex:4257-4328.
% AI-transcribed by Claude Opus 5.5 (claude-opus-5-5, adaptive thinking, effort high), Message Batches API, batch ps2 (msgbatch_01NVwr1Yuo8qJf5rykeGBXPK), 2026-11-12; not yet checked by eye.
\scanpage{134-5/83}
\mspage{270}

here are those associated to the given structure $M_c$ . The case dealt with
in section 57 (p. 143) is \xout{\unsure{the one}}\inserted{essentially} (it seems) \deleted{essentially} the one when
these notions are the ones canonically defined in terms of the category
structure of M alone. In view of the inclusions (1), we could throw in
a handfull more \inserted[typed]{obviously} equivalent conditions, using $M_{\infty}$ , $M_{\underline{W}}$ or $W_{\infty}$ , $\underline{W}_M$\inserted{instead of $M_c$ etc}, but
it seems this would confuse the picture rather than complete it .

Next thing is to look at the canonical map
\[ (5) \qquad \hand{\overline{\Gamma}}(X) \to \hand{\pi}_o(X) \]
defined for any object X of M , where $\hand{\Gamma}(X)$ denotes the set of sections
of X, and $\hand{\overline{\Gamma}}(X)$ denotes the set of \unsure{cotrresponding|corresponding} homotopy classes of sec-%
tions. (This map was considered in a \xout{\unsure{special case}} slightly different case
on page 144.) The map (5) can be viewed as the particular case of
\[ \overline{\mathrm{Hom}}(Y,X) \to \mathrm{Hom}_{\mathrm{Sets}}(\hand{\pi}_o(Y),\hand{\pi}_o(X)) , \]
obtained when $Y = e_M$ hence $\hand{\pi}_o(Y) =$ one-point set. By the standard des-%
cription
\[ (6) \qquad \hand{\pi}_o(X)\simeq \underline{\hand{\pi}}_o(C_{/X}) \quad ( = \hand{\pi}_o(i^*(X))) \]
we immediately get that (5) is always surjective, as any element in $\hand{\pi}_o(X)$
is induced by an element $x \to X$ of $C_{/X}$ , hence by a section
\[ e_M \to x \to X \quad , \]
where $e_M \to x$ is a section of the (contractible) object x in C. But (5)
is in fact \uline{bijective} . To see this, let's note that the map (5) is
isomorphic to the corresponding map\inserted{in (Sets)}, with M replaced by C\^{} and X by $i^*(X)$,
C remaining the same, as follows from \xout{the} (6) and the formula
\[ \hand{\Gamma}(X)\simeq \hand{\overline{\Gamma}}(i^*(X)) \quad . \]
The latter is a particular case of the

\uline{Lemma} 3. Let Y,X be two objects of M, with Y in C. Then the natural map
\[ (7) \qquad \overline{\mathrm{Hom}}(\over{\unsure{X}}{Y},\over{\unsure{Y}}{X}) \to \overline{\mathrm{Hom}}(i^*(Y),i^*(X)) \]
between sets of homotopy classes of maps (in M and in C\^{}respectively, the
latter being endowed with the contractibility structure generated by C)
is bijective.

The verification is tautological, due to the fact that for any
object I of C (which is going to play the role of a homotopy interval for
deformations), we got
\[ \mathrm{Hom}_M(\mathrm{IxY},X) \xrightarrow{\sim} \mathrm{Hom}_{C\hat{}}(\mathrm{IxY},i^*(X)). \]

\margin{83}
